\documentclass[a4paper,11pt]{article}
\usepackage{exam-style-842-unit}
\examinehead{842 信息化概论 · 第09单元练习册}{练习一 参考答案与解析}
\begin{document}

\unitanswerscover{842 信息化概论 · 单元练习册}{第09单元参考答案}{练习一（智慧农业与精准农业）的答案与解析；每道判断题均给出改正后的正确表述，综合题另给出答题框架、给分点与分值分配、常见失分。}

\answerpart{一、单项选择题（10 题，每题 3 分，共 30 分）}

\q B。\begin{analysis}智慧农业是以数据为关键生产要素、以人工智能为核心驱动力，依托物联网、大数据、云计算等实现农业生产全过程感知、决策、控制和服务的现代农业形态。A 劳动力投入是传统农业的主要投入；C 化肥农药是传统农业“高投入”的典型表现；D 土地规模是经营条件，不是智慧农业的核心驱动。\end{analysis}

\q B。\begin{analysis}精准农业（Precision Agriculture）的核心思想是承认农田内部存在空间差异，按地块、按管理单元“定位、定时、定量”地投入，即因地制宜、按需投入。A 大面积统一施肥正是精准农业要克服的“一刀切”；C 完全依靠人工经验是传统农业特征；D 只追求产量最高忽视效率与生态，也不符合精准农业“高产、优质、高效、低耗、环保”的目标。\end{analysis}

\q C。\begin{analysis}精准农业的一般实施过程为：信息采集（土壤、作物、环境的空间差异数据）$\rightarrow$ 差异分析（划分管理单元、生成处方图）$\rightarrow$ 决策生成（确定施肥、灌溉、喷药方案）$\rightarrow$ 变量作业（按处方精确实施）。因此第一步是信息采集；A、B、D 都发生在信息采集之后。\end{analysis}

\q B。\begin{analysis}处方图是在 GIS 中把地块差异分析结果与农艺决策落到具体空间位置的地图，是变量作业装备的作业依据。A 地形图只反映地形地貌；C 气象图反映气象要素分布；D 土地利用现状图反映土地利用类型，均不含施肥、灌溉等作业决策信息。\end{analysis}

\q B。\begin{analysis}遥感通过搭载在卫星、飞机、无人机上的传感器获取地物光谱信息，传感器可实时采集土壤墒情、温度、养分等数据，二者是获取农田空间差异信息的主要技术。A 变量施肥机是作业执行装备；C 电商平台服务流通环节；D 专家系统负责分析决策，均不承担信息获取任务。\end{analysis}

\q A。\begin{analysis}物联网的感知层通过传感器、RFID、GPS、摄像头等“随时随地获取物体信息”，承担智慧农业的“全面感知”任务。B 区块链主要解决数据可信与溯源；C 云计算提供存储与算力；D 虚拟现实用于模拟与可视化，均不承担前端的全面感知。\end{analysis}

\q D。\begin{analysis}智慧农业的支撑技术包括物联网、大数据、云计算、人工智能、3S、智能装备与区块链等。A、B、C 均在支撑技术之列；D 人畜力耕作是传统农业的生产方式，与智慧农业以智能装备代替人畜力的方向相反。\end{analysis}

\q B。\begin{analysis}三者是农业信息化不同发展阶段的形态：数字农业强调农业要素与过程数字化，是基础；精准农业在数字农业基础上强调空间差异与变量作业，是生产环节的深化；智慧农业在二者基础上强调以人工智能实现全链条“感知—决策—执行”闭环，是高级阶段。A 混淆了三个概念；C 把包含关系说反了；D 数字农业恰恰以信息技术为支撑。\end{analysis}

\q B。\begin{analysis}我国智慧农业近年来发展较快，已建成一批示范基地和无人农场，但总体上仍处于起步与示范推广阶段：感知设备、智能农机成本较高，数据标准和共享机制不健全，小农户应用门槛高。A、C 高估了普及程度；D 低估了已取得的进展。\end{analysis}

\q B。\begin{analysis}精准农业依据地块差异生成处方图并按处方变量投入，避免了传统“大水大肥”造成的过量施肥、过量用药，因而能在稳定产量的同时减少化肥农药用量。A 减少种植面积、C 降低产量目标都不是精准农业的出发点；D 用人工替代机械与事实相反。\end{analysis}

\answerpart{二、多项选择题（4 题，每题 4 分，共 16 分）}

\q ABC。\begin{analysis}智慧农业的基本特征可概括为感知化（全面感知）、数字化（数据驱动）、网络化（互联互通）、智能化（智能决策与控制）、精准化（按需投入）和全链条化。D 错误：粗放化与智慧农业“精准、高效、低耗”的方向完全相反，是传统农业的特征。\end{analysis}

\q ABC。\begin{analysis}精准农业的主要支撑技术包括遥感与传感器（信息获取）、地理信息系统（空间分析与处方图生成）、全球定位系统（定位与导航）以及变量作业装备等。D 错误：精准农业正是要克服“凭人工经验估算、统一投入”的做法，经验估算不是其支撑技术。\end{analysis}

\q ABCD。\begin{analysis}精准农业的实施过程包括信息采集、差异分析、决策生成（处方）和变量作业四个环节，四者依次衔接、缺一不可，作业后还要进行效果评价与反馈。四项全对。\end{analysis}

\q ABC。\begin{analysis}我国发展智慧农业的主要问题：农村信息基础设施仍较薄弱、高端传感器与智能装备依赖进口、数据标准与共享机制不健全、建设与运维成本高、农民数字技能不足、小规模经营与服务组织不健全等。D 错误：我国农业劳动力总体呈老龄化、兼业化和结构性短缺态势，并非“严重过剩”。\end{analysis}

\answerpart{三、判断题（6 题，每题 2 分，共 12 分）}

\q ×。改正：智慧农业是以数据和人工智能为核心、实现农业生产全过程“感知—决策—执行”闭环的全链条智能化，远程视频监控只是其中感知环节的一种具体应用，不能等同于智慧农业。\begin{analysis}把智慧农业等同于“装摄像头、手机看画面”，是只看到感知环节、丢掉了数据分析和智能决策，属于典型的概念窄化。\end{analysis}

\q √。\begin{analysis}精准农业的核心正是根据农田空间差异，按管理单元实施定位、定时、定量的变量投入，从而避免“一刀切”。\end{analysis}

\q ×。改正：精准农业信息采集后必须经过差异分析、生成处方图（决策生成），才能指导变量作业；缺少差异分析与决策，变量作业就失去了依据。\begin{analysis}精准农业四环节是“信息采集 $\rightarrow$ 差异分析 $\rightarrow$ 决策生成 $\rightarrow$ 变量作业”，顺序不能跳步。\end{analysis}

\q ×。改正：数字农业、精准农业和智慧农业是农业信息化不同发展阶段的形态，依次递进、各有侧重，不是同一概念的不同叫法。\begin{analysis}数字农业强调数字化基础，精准农业强调空间差异与变量作业，智慧农业强调人工智能驱动的全链条智能，三者内涵不同。\end{analysis}

\q √。\begin{analysis}智慧农业是在数字农业、精准农业基础上，以物联网、大数据、人工智能等为支撑发展起来的高级阶段形态，体现为全链条的自主感知、智能决策与精准控制。\end{analysis}

\q √。\begin{analysis}变量作业装备的作业量由处方图实时控制，而处方图来自地块差异分析和农艺决策，因此地块差异信息与处方图是变量作业的前提。\end{analysis}

\answerpart{四、填空题（6 空，每空 1 分，共 6 分）}

\q （1）智能化。\begin{analysis}智慧农业的技术特征通常概括为感知化、数字化、网络化、智能化和精准化，其中智能化是区别于数字农业、精准农业的关键特征。\end{analysis}

\q （2）差异分析。\begin{analysis}精准农业实施过程为：信息采集 $\rightarrow$ 差异分析 $\rightarrow$ 决策生成 $\rightarrow$ 变量作业，差异分析的结果是管理单元划分与处方图。\end{analysis}

\q （3）按需投入（或变量投入）。\begin{analysis}精准农业的核心思想是“因地制宜、按需投入”，即根据地块空间差异决定投入的种类、时间和数量。\end{analysis}

\q （4）变量作业。\begin{analysis}变量施肥机、变量喷药机、智能灌溉设备均根据处方图实时调整作业量，属于变量作业技术。\end{analysis}

\q （5）示范推广。\begin{analysis}我国智慧农业已建成一批示范基地和无人农场，但总体上处于起步和示范推广阶段，尚未全面普及。\end{analysis}

\q （6）数据。\begin{analysis}智慧农业以数据为关键生产要素，以人工智能为核心驱动，通过数据采集、汇聚、分析形成决策。\end{analysis}

\answerpart{五、简答题（4 题，每题 6 分，共 24 分）}

\q 参考答案要点（每点 2 分）：
（1）概念：智慧农业是以物联网、大数据、云计算、人工智能等现代信息技术为支撑，以数据为关键生产要素，实现农业生产全过程信息感知、定量决策、智能控制、精准投入和个性化服务的现代农业形态。
（2）基本特征：感知化（全面感知）、数字化（数据驱动）、网络化（互联互通）、智能化（智能决策与控制）、精准化（按需投入）、全链条化（覆盖产前、产中、产后）。
（3）主要支撑技术：物联网（感知与传输）、大数据与云计算（存储与计算）、人工智能（分析与决策）、3S 技术（定位、遥感与空间分析）、智能装备（无人机、智能农机）、区块链（溯源与信任）。

\q 参考答案要点（每点 1.5 分）：
（1）概念：精准农业是基于农田空间差异，综合运用 3S 技术、传感器和变量作业技术，实现定位、定时、定量按需投入的农业生产方式，目标是高产、优质、高效、低耗、环保。
（2）信息采集：利用遥感、GPS／北斗、传感器、无人机采集土壤养分、墒情、作物长势、病虫害等空间差异数据。
（3）差异分析：在 GIS 中处理数据、划分管理单元，生成反映地块差异的图层与处方图。
（4）决策生成：依据差异分析与农艺模型、专家系统，确定各地块施肥、灌溉、喷药的种类与用量。
（5）变量作业：变量施肥机、变量喷药机、智能灌溉与自动驾驶农机按处方图精确作业，并评价效果、反馈修正。

\q 参考答案要点（每类 1.5 分）：
（1）信息获取类：遥感、GPS／北斗、各类传感器、无人机，负责采集空间差异数据。
（2）信息处理类：GIS 空间分析、决策支持系统、农业专家系统，负责差异分析、处方生成与决策。
（3）变量作业类：变量施肥机、变量喷药机、智能灌溉设备、自动驾驶农机，负责按处方精确作业。
（4）支撑保障类：农村网络与数据中心、数据标准与共享平台、农业社会化服务组织，负责数据流通与落地服务。
（说明：按技术功能划分，前三类构成“获取—处理—作业”主干。）

\q 参考答案要点：
（1）现状（2 分）：农村信息基础设施明显改善，农业物联网、大数据平台、智能农机和无人农场已在示范基地应用，农产品电商与农业信息服务快速发展，国家出台数字乡村等政策支持。
（2）问题（4 分，答出 4 点即可）：一是农村信息基础设施仍较薄弱，存在“最后一公里”问题；二是高端传感器、芯片与智能装备依赖进口，核心技术与装备受制于人；三是数据标准与共享机制不健全，数据孤岛现象突出；四是建设与运维成本高，小农户经营规模小、应用门槛高；五是农民数字技能与复合型人才不足；六是农业信息服务体系与商业模式尚不健全、区域发展不平衡。

\answerpart{六、综合题（1 题，共 12 分）}

\q 参考答案与答题框架：

\textbf{【破题句】}该示范区是以物联网全面感知为基础、以数据平台和智能决策为中枢、以智能装备为执行手段的“感知—决策—执行”闭环智慧农业雏形，其中精准施肥、精准灌溉集中体现了精准农业“信息采集—差异分析—决策生成—变量作业”的实施过程。

\textbf{（1）支撑技术与任务（4 分）}
\begin{analysis}
答出“技术—任务”对应关系，每对 1 分，答出 4 对即可：
（1）物联网感知技术（墒情传感器、气象站）：全面感知土壤水分、养分和农田小气候，是数据源头；
（2）无人机遥感：低空获取作物长势、病虫害和受灾情况，弥补地面传感器点位有限的不足；
（3）大数据平台与云计算：汇聚、存储、处理多源数据，形成县域农业数据底座；
（4）人工智能与模型／专家系统：分析数据、生成施肥灌溉与农事建议，是“决策中枢”；
（5）智能农机（自动驾驶拖拉机、变量施肥机）：按处方执行精准作业，是“执行末端”；
（6）手机端应用：把决策结果转化为农户可操作的信息服务，打通“最后一公里”。
\end{analysis}

\textbf{（2）按需施肥、按需灌溉的实施过程（5 分）}
\begin{analysis}
按四环节给分：信息采集（1 分）$\rightarrow$ 差异分析（2 分，须提到划分管理单元、生成处方图）$\rightarrow$ 决策生成（1 分）$\rightarrow$ 变量作业（1 分），结合水稻、柑橘举例可酌情加印象分但不单独计分。答题要点：
（1）信息采集：传感器采集土壤墒情、氮磷钾含量，无人机获取长势，形成地块级数据；
（2）差异分析：在 GIS 中把地块划分为若干管理单元，生成土壤养分差异图与灌溉需水差异图（处方图）；
（3）决策生成：依据作物需肥需水规律和模型，确定每个单元的施肥量、施肥时期与灌水量、灌溉时机；
（4）变量作业：变量施肥机按处方图实时调节施肥量，智能灌溉系统按墒情阈值自动启闭，自动驾驶农机保证行走精度；
（5）效果评价：对比产量、品质与投入成本，反馈修正模型与处方。
举例：柑橘园按树龄、地势和土壤养分分区施氮，避免坡地过量流失；水稻田按墒情传感器阈值控制灌溉，实现浅湿交替、节水增产。
\end{analysis}

\textbf{（3）小农户地区推广的困难（3 分）}
\begin{analysis}
答出 3 点、每点 1 分，言之成理即可：
（1）经营规模小、地块细碎，智能装备作业效率与投资回收期不划算，规模不经济；
（2）传感器、智能农机与平台建设运维成本高，小农户单独承担不起；
（3）农民年龄偏大、数字技能不足，对手机端建议的接受度和使用能力有限；
（4）数据标准与共享机制不健全，平台之间数据孤岛，服务难以持续；
（5）网络、电力等基础设施和售后维修服务不完善，设备“用不好、修不了”。
对策可简述（不单独计分）：发展土地流转与托管、培育社会化服务组织、政府补贴与示范基地培训、降低设备成本、健全数据标准。
\end{analysis}

\textbf{【给分点与分值分配】}（1）4 分：技术—任务对应 4 对；（2）5 分：信息采集 1 分、差异分析 2 分、决策生成 1 分、变量作业 1 分；（3）3 分：困难 3 点。

\textbf{【常见失分】}一是只堆“物联网、大数据、人工智能”名词，不说明各自承担的任务，技术与环节脱节；二是把精准农业四环节顺序写乱或漏掉差异分析，直接由采集跳到作业；三是只讲困难不提任何对策，或把困难写成“农民素质低”一句空话，缺少规模、成本、数据、服务等具体维度；四是全篇不结合水稻、柑橘等生产实际；五是把“智慧农业”等同于“远程监控”，概念把握不准。

\end{document}
